\documentstyle[%


righttag%


,11pt%


,amssymb%


,russian%               remove in Eng.


,izvru%                 Set izven% in Eng.


]{amsart}





% 


\numberwithin{equation}{section}


\newcommand{\Mod}{\operatorname{mod}}


\newcommand{\re}{\operatorname{Re}}


\newcommand{\card}{\operatorname{card}}


\newcommand{\tts}[1]{}





\theoremstyle{plain}


\theorembodyfont{\it}


\newtheorem{lems}{\hskip\parindent Лемма}[section]





%\hoffset=-15mm%                  For Eng.


\setcounter{page}{3}


\god{2000}


\nomer{8~(459)} %                        Remove in Eng.


%\nomer{Vol.\,44, No.\,8}%               Set in Eng.


%\str{\pageref{begin}--\pageref{end}}%  Set in Eng.


\udk{511.216}





\author{\it И.А.\,Аллаков}


% NB:   ^^ remove in Eng.


\title{О представлении четных чисел суммой двух нечетных простых


чисел из арифметической прогрессии}





\date{}


%\pagestyle{myheadings}%         Set in Eng.


\pagestyle{plain} %              Remove in Eng.


\begin{document}


%\thispagestyle{plain}%          Set in Eng.


\maketitle


\label{begin}





\section{Введение}





Пусть $X$ и $P$ --- достаточно большие вещественные числа, а


$N$ --- натуральное число с условием $\sqrt{N}<X\le N$;


$p$, $p_1$, $p_2$ ---


простые числа; $D=p^{\nu}$, где $p>2$, $\nu$ --- положительное целое;


$M_D(X)$ --- множество четных натуральных чисел $n\le X$, которые (возможно)


не представимы в виде


\begin{equation}


n=p_1+p_2,\quad p_i\equiv l_i(\Mod D),\quad (l_i, D)=1,\quad i=1,2;


\end{equation}     


$E_D(X)=\card M_D(X)$; $R(n)$ --- число представлений $n$ в виде (1.1);


$c$, $c_j$ ($j=1,2,\dotsc$) --- некоторые положительные постоянные,


$\ll$ --- символ Виноградова, $\varphi (a)$ --- функция Эйлера.





В [1] была получена асимптотическая формула для $R(n)$, справедливая для всех


четных $n\le X$, за исключением не более $E_D(X)\ll X\ln^{-A}X$ (где $A>0$


--- произвольная постоянная) значений $n$ из них. Затем в [2], [3] для


$E_D(X)$ в случае $D=1$ соответственно доказано,


что $E_1(X)<X\exp(-c\sqrt{\ln X})$


и $E_1(X)<X^{1-\delta}$, $0<\delta<1$. В [4] получена оценка снизу $R(n)$


при $D=1$ для всех $n\le X$, за исключением не более $X^{1-\delta}$ значений


$n$ из них.





В данной работе соединением методик работ [1], [2] доказана





\begin{thm}


При $D=p^{\nu}$ и $D\ll \ln^A X$ справедливы оценки


$$


E_D(X)\ll X\varphi^{-1}(D)\exp(-c_1\sqrt{\ln X})


$$


и для $n\notin M_D(X)$, $n\le X$


$$


R(n)\gg\frac{n}{\varphi(D)\ln^2 n}


\bigg(1-\frac{\ln^A n}{\exp(c_2\sqrt{\ln n})}\bigg)


\exp\bigg(-\frac{c_2}{4}\sqrt{\ln n}\bigg), 


$$


где постоянные $c_1$ и $c_2$ не зависят от $A$.


\end{thm}





Пусть $s$ --- комплексная переменная, $\chi_q(n)$ --- характер Дирихле


по модулю $q$ и $L(s,\chi_q)$ --- $L$-функция Дирихле, определяемая при


$\re s>1$ равенством


$L(s,\chi_q)=\sum\limits_{n=1}^\infty\chi_q(n)n^{-s}$. Если 


не существует вещественный (исключительный) нуль $\beta$ с условием


$\beta>1-c_3\ln^{-1}q$\; $L$-функции Дирихле для всех $q\le P$


(в этом случае положим $E=0$), то методикой из данной работы можно получить 


асимптотическую формулу для $R(n)$.


В случае существования такого исключительного нуля


(тогда положим $E=1$) также получается асимптотическая формула, но


в этой формуле вместе с обычным главным членом будет участвовать член,


соответствующий исключительному нулю (порядок которого по сути одинаков


с главным членом), т.\,е. справедлива





\begin{thm}


Если $D=p^\nu$ и $D\ll \ln^A X$, то для всех $n\le X$, исключая


не более чем\linebreak $c_4 X\varphi^{-1}(D)\exp(-c_1\sqrt{\ln X})$ 


значений из них, справедлива формула


$$


R(n)=\frac{J(n)\sigma(n)}{\varphi^2(D)\ln^2 n}+


E\frac{J_{\beta}(n)\wt\sigma(n)}{\varphi^2(D)\ln^2 n}+


O\bigg(\frac{n(\varphi(D)\ln^2 n)^{-1}}{\exp(c_5\sqrt{\ln n})}\bigg),


$$


где


\begin{gather*}


J(n)=\sum\begin{Sb} n=n_1+n_2\\ n_1,n_2\ge 2\end{Sb}1,\quad


J_{\beta}(n)=\sum\begin{Sb}n=n_1+n_2\\ n_1,n_2\ge 2\end{Sb}


(n_1n_2)^{(\beta -1)},\\


%


\sigma(n)=2D\prod_{p>2}\frac{p(p-2)}{(p-1)^2}\,


\prod\begin{Sb}p\setminus D\\ p>2\end{Sb}\frac{(p-1)^2}{p(p-2)}\,


\prod\begin{Sb}p\setminus n\\


p\not\setminus D,\ p>2\end{Sb}\frac{p-1}{p-2}, \\


%


\wt\sigma(n)=\prod\begin{Sb}p \setminus rd\\


p\not\setminus n\end{Sb}\frac1{p-1}\,


\prod\begin{Sb}p\not\setminus nrd\\ p>2\end{Sb}\frac{p(p-2)}{(p-1)^2}\,


\prod\begin{Sb}p\setminus ndr\\ p>2\end{Sb}\frac{p}{p-1}, 


\end{gather*}


где $d=(q,D)$, $r$ --- модуль вещественного (исключительного) характера


$\wt \chi_r$ и $\beta$ --- исключительный нуль функции $L(s,\wt \chi_r)$.


\end{thm}





Заметим, что методика работы [3] не позволяет получить даже такую 


асимптотическую формулу. Отметим также, что в


[5] приводится оценка $E_D(X)\ll \varphi^{-1}(D)X^{1-\delta}$, однако она 


доказана другим путем и только в случае, когда $D=p$ --- простое число.





Известно, что изучение функции $R(n)$ можно связать с изучением функции


$\prod_n(X,D)$, означающей число пар простых чисел $p$,


$p+2n$ из интервала


$(0;X)$, принадлежащих соответственно арифметическим прогрессиям 


$Dt_1+l_1$, $Dt_2+l_2$ с условием $1\le l_1,l_2\le D$,


$(l_1,D)=1,\,(l_2,D)=1$ [1].





Методика, используемая в данной работе, дает возможность доказать аналогичный


результат и для $\prod_n(X,D)$, а именно, справедлива





\begin{thm}


Если $D=p^\nu$ и $D\ll \ln^A X$, то для каждого целого


$0< 2n\le X \ln^{-A}X$, $2n\equiv l_1 - l_2(\Mod D)$, исключая не более чем 


$c_4 X\varphi^{-1}(D)\exp(-c_1\sqrt{\ln X})$


значений из них, справедлива оценка


$$


\prod\nolimits_n(X,\,D)\gg


\frac{n}{\varphi(D)\ln^2 n}\bigg(1-\frac{\ln^A n}{\exp


(c_2\sqrt{\ln n})}\bigg)\exp\bigg(-\frac{c_2}4\sqrt{\ln n}\bigg),


$$ 


где постоянные $c_1$ и $c_2$ не зависят от $A>0$.


\end{thm} 





Полученные результаты являются усилением результатов [1] и обобщением


результата [2].





Доказательство теоремы 3 в основном аналогично доказательству теоремы 1, причем


необходимые изменения для вывода теоремы 3 можно посмотреть в [1]. В ходе 


доказательства теоремы 1 сначала установим справедливость теоремы 2.





\section{Основные леммы}





Пусть $\Lambda (n)$ --- функция Манголдта, определяемая равенством


$$


\Lambda (n)=


\begin{cases}


\ln n,&\text{если } n=p^k;\\


0&\text{в остальных случаях}.


\end{cases}


$$ 


Обозначим $\psi(x,\chi_m)=\sum\limits_{n\le x}\Lambda(n)\chi_m (n)$ и


$$


\delta_\chi=


\begin{cases}


1,&\chi_m=\chi_m^\circ\text{ (главный характер)};\\


0,&\chi_m\ne\chi_m^\circ.


\end{cases}


$$


Символы $\sum\limits_{\chi_q}$ и


$\sideset{}{'}\sum\limits_{a=1}^q$ обозначают


соответственно суммирование по всем характерам $\Mod q$ и по приведенной 


системе вычетов $\Mod q$.





\begin{lems}[{\series{m}\shape{n}\selectfont [6], \S\,19; [2], п.\,4}]


Для достаточно больших $X$ и $N$


с условием $3\le n\le X$, $\sqrt{N}<X\le N$ и для всех 


$m\le\exp(c_5\sqrt{\ln N})$ справедливы следующие равенства$:$


\begin{itemize}


\item[a)] если $E=0$ и $E=1$ и $m$ --- модуль характера $\chi_m$ не делится


на $r$ --- ведущий модуль исключительного характера $\wt\chi_m$, то 


$\psi(n, \chi_m)=\delta_{\chi}n+\rho_n$;


\item[b)] если же $E=1$ и $r\setminus m$, то 


$\psi(n, \chi_m)=\delta_{\chi}n-\beta^{-1}n^\beta+\rho_n$,


причем во всех случаях для $\rho_n$ справедлива оценка


$$


\rho_n\ll X\exp(-c_6\sqrt{\ln N}).


$$


\end{itemize}


\end{lems}





\begin{lems}[{\series{m}\shape{n}\selectfont [7]}]


Если $\chi'(t)$ и $\chi''(t)$ --- характеры Дирихле 


соответственно по $\Mod Q'$ и $Q''$, то произведение $\chi'(t)\chi''(t)$ ---


главный характер тогда и только тогда, когда $\chi'(t)=\ov \chi''(t)$


для всех $(t,Q)=1$ и $\chi'$ --- производный характер, порожденный одним из


характеров по $\Mod d$, где $d=(Q',Q'')$, $Q=Q'Q''d^{-1}$ и $\ov\chi''(t)$


--- характер, сопряженный с $\chi''(t)$.


\end{lems}





\begin{lems}[{\series{m}\shape{n}\selectfont [7]}]


Для $(j,k)=1$, $d\setminus k$ и $(h,d)=1$ имеют 


место соотношения


$$


\sideset{}{'}\sum^k\begin{Sb}l=1 \\ l\equiv h(\Mod d)\end{Sb}


e\bigg(\frac{jl}{k}\bigg)=


\begin{cases}


0,& \text{если } (d, kd^{-1})>1;\\


\mu \left(\frac{k}{d}\right)e\left(\frac{j\theta h}{d}\right),& \text{если }


(d, kd^{-1})=1,


\end{cases}


$$


где $e(\alpha) =e^{2\pi i\alpha}$, $\mu (k)$ --- функция Мебиуса и $\theta$


--- решение сравнения $kd^{-1}z\equiv 1(\Mod d)$.


\end{lems}





\section{Основные обозначения и деление единичного интервала}





Пусть 


$$


P_1=\exp\bigg(c_6\frac{\sqrt{\ln N}}{200}\bigg),\quad


P_2=P_1^{1/4}.


$$





Если $E=0$ или же $E=1$ и $P_2<r$, то положим $P=P_2$, в остальных случаях


$P=P_1$. Полагая


$Q=NP^{-1}$,


сегмент $[Q^{-1}, 1+Q^{-1}]$ делим на основные и дополнительные интервалы.


При $a\le q\le P$ через $M(q, a)$ обозначим закрытый интервал


$[aq^{-1}-(qQ)^{-1}, aq^{-1}+(qQ)^{-1}]$, $(a,q)=1$. Ясно, что основные


интервалы не пересекаются и $M(q, a)\subset [Q^{-1}, 1+Q^{-1}]$.





Через $T$ обозначим множество тех точек $\alpha$, $Q^{-1}<\alpha<1+Q^{-1}$,


которые не содержатся ни в каком $M(q, a)$. В дальнейшем объединение всех 


$M(q, a)$ назовем большой дугой, а $T$ --- малой дугой.





Введем функции


\begin{gather}


S_i(X,\alpha)=\sum_{\chi_{D}}\ov\chi_{D}(l_i)\sum_{2<p_i\le X}


\chi_D(p_i)\ln p_i e(p_i\alpha),\quad i=1, 2;\\


%


\mathrm g_u^{(i)}(x,\alpha)=\sum_{2<n\le X}n_i^{u-1}e(n_i\alpha),


\quad i=1, 2;\\


%


V_i(X,\alpha,q,a)=R(q)\frac{\mu(qd^{-1})}{\varphi(qd^{-1})}e


\bigg(\frac aq N_1l_i\bigg)\mathrm g_1^{(i)}


(X,\eta),\quad i=1, 2,


\end{gather} 


где $\alpha=aq^{-1}+\eta$, $d=(q, D)$, $N_1=\mathrm g qd^{-1}(\Mod q)$


($N_1$ --- наименьший положительный вычет числа


$\mathrm g qd^{-1}$ по $\Mod q$), $\mathrm g$ по


$\Mod d$ определяется из $\mathrm g qd^{-1}\equiv 1(\Mod d)$; $R(q)=1$, если


$( qd^{-1},D)=1$, и $R(q)=0$ в противном случае, обозначим


\begin{equation}


\begin{array}{c}


S=S(X,\alpha)=S_1(X,\alpha) S_2(X,\alpha),\\[2mm]


V=V(X,\alpha,q,a)=V_1(X,\alpha,q,a) V_2(X,\alpha,q,a).


\end{array}


\end{equation} 


Тогда имеем


\begin{equation}


S(X,\alpha)=\varphi^2(D)\sum_{2<n\le 2X}R(X,n) e(\alpha n),


\end{equation}


где


\begin{equation}


R(X,n)=\sum\begin{Sb}n=p_1+p_2\\ 2<p_1, p_2\le X\\


p_1\equiv l_1,\ p_2\equiv l_2 (\Mod D) \end{Sb}


\ln p_1\ln p_2,


\end{equation}


и


\begin{equation}


V=R(q)\frac{\mu^2(qd^{-1})}{\varphi^2 (qd^{-1})}


\sum_{2<n\le 2N}J(N,n)e\bigg(\frac aq(N_1(l_1+l_2)-n)\bigg)e(\alpha n),


\end{equation}


где


\begin{equation}


J(N,n)=\sum\begin{Sb}n=n_1+n_2\\ 2\le n_1, n_2\le N\end{Sb}1.


\end{equation}


Очевидно,


\begin{equation}


n/2<J(N,n)\ll N\quad \left(J(N,n)=2(n-2)\right),


\end{equation}


если $2<n\le N$.


Если $1/2\le u\le 1$, то суммирование по частям дает ([8], лемма 3.5)


\begin{equation}


\mathrm g_u^{(i)}(X,\alpha)\ll\min(\|\alpha\|^{-u},X^u),


\end{equation}


где $\|\alpha\|$ --- расстояние от $\alpha$ до ближайшего целого числа.





\section{Малые дуги}





\begin{lems}


При достаточно большом $N$ справедливы оценки


\begin{align}


\int_T |S(N,\alpha)|^2 d\alpha&\ll N^2D^2\varphi(D)P^{-1}\ln^{12}N,\\


%


\int_T\bigg|\sum_{q\le P}\sideset{}{'}\sum_{a=1}^q


V(N,\alpha)\bigg|^2 d\alpha&\ll N^3P^{-2}(D\ln\ln P)^4.


\end{align}


\end{lems}





\begin{pf}


В силу (3.4) имеем


\begin{equation}


\int_T|S(N,\alpha)|^2 d\alpha\ll\max\limits_{\alpha\in T}|S_1(N,\alpha)|^2


\int_T|S_2(N,\alpha)|^2 d\alpha.


\end{equation}


Здесь


$$


\int_T|S_2(N,\alpha)|^2 d\alpha\le\int_{Q^{-1}}^{1+Q^{-1}}


|S_2(N,\alpha)|^2 d\alpha=\varphi^2(D)


\sum\begin{Sb}2<p\le N\\ p_2\equiv l_2(\Mod D)\end{Sb}


\ln^2p_2.


$$


В силу теоремы 5.2.1 из [9] при $N\ge 2$ и $D\le N^{1/2}$


$$


\sum_{p\le N,\ p\equiv l(mod\,D)}\ln p\ll N\varphi^{-1}(D).


$$


Поэтому 


\begin{equation}


\int_T|S_2(N,\alpha)|^2 d\alpha\ll \varphi(D)N\ln N.


\end{equation}


Для оценки $\max\limits_{\alpha\in T}|S_1(N,\alpha)|^2$ используем


результат работы [10]: если


$R<q\le NR^{-1}$, $1\le R\le N^{1/3}$, $(a,q)=1$,


$|\alpha -aq^{-1}|\le 2R(qN)^{-1}$, то 


\begin{equation}


S_i (N,\alpha)\ll NR^{-1/2}D(\ln N)^{11/2}.


\end{equation}


Согласно теореме Дирихле существуют такие $q\le Q$ и $\alpha$ с условием


$1\le a\le q$, $(a, q)=1$, для которых $|\alpha -aq^{-1}|<(qQ)^{-1}$. Это


означает, что $\alpha\in M(q,a)$, если $q\le P$. Значит, для $\alpha\in T$


имеем $q>P$ и, следовательно, в (4.5) можем полагать $R=P$. Теперь из


(4.3)--(4.5) следует (4.1).





Оценка (4.2) доказывается так же, как оценка (5.9) из [2].


\end{pf}





\section{Главные дуги}





A. Случай $P=P_2$. Из (3.7) следует, что


\begin{equation}


\sum_{q>y}\,\sideset{}{'}\sum_{a=1}^q V(N,\alpha,q,a)=\sum_{2<n\le 2N}


J(N,n)\sigma(y,n)e(n\alpha),


\end{equation} 


где


\begin{equation}


\sigma(y,n)=\sum_{q>y}R(q)\frac{\mu^2(qd^{-1})}{\varphi^2 (qd^{-1})}


\sideset{}{'}\sum_{a=1}^q e\bigg(\frac aq(N_1(l_1+l_2)-n)\bigg).


\end{equation}





\begin{lems}


Если $\alpha\in M(q,a)$, то


$$


\sum_{q\le P}\,\sideset{}{'}\sum_{a=1}^q\int_{M(q,a)}


\bigg|\sum\begin{Sb}k\le P\\ k\ne q\end{Sb}


\sideset{}{'}\sum\begin{Sb}h=1\\ h\ne a\end{Sb}^k


V(N,\alpha, k, h)\bigg|^2


d\alpha\ll\frac{P^9}{Q}(D\ln\ln P)^4.


$$


\end{lems}





{\bf Доказательство.} Предположим, что 


$a\le q\le P$, $(a, q)=1$,


$h\le k\le P$, $(h, k)=1$, $k\ne q$, $h\ne a$ и $\alpha\in M(q, a)$.


Тогда $|\alpha-aq^{-1}|\le(qQ)^{-1}$. Следовательно,


$\|\alpha-hk^{-1}\|=\|\alpha - hk^{-1}+ aq^{-1} -aq^{-1}\|


\ge \|aq^{-1} - hk^{-1}\| - |\alpha- aq^{-1}|\ge (qk)^{-1} -


(qQ)^{-1} > (qk)^{-1}$.


Поэтому, учитывая $n\varphi^{-1}(n)\ll \ln\ln n$ (при $n\ge 3$) и


$d_1=(k,D)\le D$, из (3.3), (3.4), (3.10) получим


\begin{multline*}


\bigg|\sum\begin{Sb}k\le P\\ k\ne q\end{Sb}


\sideset{}{'}\sum^k\begin{Sb}h=1\\ h\ne a\end{Sb}


V(N,\alpha, k, h)\bigg|^2\ll


\bigg|\sum\begin{Sb}k\le P\\ k\ne q\end{Sb}


\frac{\mu^2(kd_1^{-1})}{\varphi^2 (kd_1^{-1})} 


\|\alpha -hk^{-1}\|^{-2}\varphi(k)\bigg|^2 \ll \\


%


\ll (q D\ln\ln P)^4\bigg(\sum_{k\le P}\varphi(k)\bigg)^2\ll


(q D P\ln\ln P)^4.


\end{multline*}





Таким образом,


\begin{multline*}


\sum_{q\le P}\,\sideset{}{'}\sum_{a=1}^q\int_{M(q,a)}


\bigg|\sum\begin{Sb}k\le P\\ k\ne q\end{Sb}


\sideset{}{'}\sum\begin{Sb}h=1\\ h\ne a\end{Sb}^k


V(N,\alpha, k, h)\bigg|^2 d\alpha \ll \\


%


\ll\sum_{q\le P}\,\sideset{}{'}\sum_{a=1}^q (PDq \ln\ln P)^4(qQ)^{-1}\ll


P^9 Q^{-1}(D\ln\ln P)^4.\quad\square


\end{multline*}





\begin{lems}[{\series{m}\shape{n}\selectfont [1], \S\,7}]


Для суммы $\sigma (P, n)$, 


определяемой равенством $(5.2)$, справедлива оценка


$$


\sigma (P, n)\ll D^2 P^{-1}\tau(n)(\ln\ln N)^3, 


$$


где $\tau(n)$ --- число натуральных делителей $n$.


\end{lems}





Из леммы 5.2 и из (5.1), (3.9) следует 





\begin{lems}


При достаточно большом $N$ имеет место оценка


$$


\int_{Q^{-1}}^{1+Q^{-1}}


\bigg|\sum_{q> P}\,\sideset{}{'}\sum_{a=1}^q V(N,\alpha,k,h)\bigg|^2


d\alpha \ll N^3P^{-2}(D \ln N)^4.


$$


\end{lems}


 


\begin{lems}


Если $a\le q\le P$, $(a, q)=1$ и $\alpha\in M(q, a)$, то


$$


S_i(N,\alpha)-V_i(N, \alpha, q, a)\ll 


N^2q^{-1}Q^{-1}\exp (-c_6\sqrt{\ln N})^4.


$$


\end{lems}





\begin{pf}


Ясно, что


\begin{equation}


\bigg|\sum_{p_i\le N}\chi_{D} (p_i)\ln p_i e(\alpha p_i) -


\sum_{p_i\le N,\ p_i\not\setminus q}\chi_{D} (p_i)\ln p_i e(\alpha p_i)


\bigg|\le\ln q.


\end{equation}


Используя свойство ортогональности характеров, вычитаемую сумму в левой части


(5.3) можно написать в виде


\begin{equation}


\frac1{\varphi(q)}


\sum_{\chi_q}\bigg(\sum_{j=1}^q {\ov \chi_q}(j) e\bigg(\frac{a j}{q}\bigg)


\bigg)\sum_{p_i\le N}(\ln p_i)\chi_{D}(p_i)\chi_q(p_i) e(p_i\eta).


\end{equation}





Обозначим


\begin{equation}


A_i=A_i(\chi_{D}, \chi_q, q, a)=


\frac1{\varphi(q)}


\sum_{\chi_D, \chi_q}\,\sum_{j=1}^q {\ov \chi_D}(l_i){\ov \chi_q}(j) e


\bigg(\frac{a j}{q}\bigg)


\end{equation}


и


\begin{equation}


G_i=G_i(\chi_D, \chi_q, N)=


\sum_{p_i\le N}(\ln p_i)\chi_D(p_i)\chi_q(p_i) e(p_i\eta),\quad i=1,2.


\end{equation}


Тогда, используя (5.3)--(5.6), из (3.1) находим


\begin{equation}


S_i(N,\alpha) = A_iG_i +O(\varphi(D)\ln q),\quad i=1,2.


\end{equation}





Сначала оценим $G_i$. Положим $\chi_D(p_i)\chi_q(p_i)= \chi_m(p_i)$,


где $m=q D d^{-1}$, тогда из (5.6) получим


\begin{multline*}


G_i=\sum_{p_i\le N}\chi_m(p_i)\ln p_i e(p_i\eta)=


\sum_{n_i\le N}\chi_m(n_i)\Lambda (n_i) e(n_i\eta)-


\sum_{


p^\nu\Big\{ \begin{array}{l}\hspace{-1mm}\scriptstyle \le N\\


\hspace{-1mm}\scriptstyle\ge 2\end{array}


}


\chi_m(p_i^\nu)\ln p_i e(p_i^\nu\eta)=\\


%


=\sum_{n_i\le N}(\psi(\chi_m, n_i)-


\psi(\chi_m, n_i-1)) e(n_i,\eta)+O(N^{1/2}).


\end{multline*}


Отсюда, используя лемму 2.1 а), находим


\begin{multline*}


G_i(\chi_m, N)-\delta_\chi\sum_{2< n_i\le N} e(n_i\eta)\ll\bigg|


\sum_{2< n_i\le N}(\rho_{n_i}-\rho_{n_i-1}) e(n_i\eta)\bigg|+\sqrt{N}\ll\\


%


\ll \bigg|


\sum_{2< n_i\le N}(\rho_{n_i}-\rho_{n_i-1})\bigg|+\bigg|\sum_{2<n_i\le N-1}


( e(n_i\eta)-e((n_i+1)\eta))


\sum_{2< t\le n_i}(\rho_t-\rho_{t-1})\bigg|+ \\


%


+\sqrt{N}\ll


|\rho_N|+ \sum_{2< n_i\le N}|


e(n_i\eta)-e((n_i+1)\eta)|\,|\rho_{n_i}|+\sqrt{N}.


\end{multline*}





Далее, принимая во внимание


$|e(n_i\eta)-e((n_i+1)\eta)|\ll|\eta|\le (qQ)^{-1}$,


будем иметь


\begin{multline}


G_i(\chi_m, N)-\delta_\chi\sum_{2< n_i\le N} e(n_i\eta)\ll |\rho_{N}|+


N(qQ)^{-1}\max\limits_{n_i}|\rho_{N}|+\sqrt{N}\ll \\


%


\ll N(qQ)^{-1}|\rho_N|\ll N^2(qQ)^{-1}\exp(-c_6\sqrt{\ln N}).


\end{multline}


Таким образом, из (5.7) и (5.8) находим


\begin{equation}


S_i(N,\alpha)=\delta_\chi A_i\sum_{2< n_i\le N} e(n_i\eta)+


O\bigg(\frac{N^2}{qQ}|A_i|\exp (-c_6\sqrt{\ln N})\bigg).


\end{equation}


Теперь найдем $A_i$. В силу леммы 2.2\;


$\chi_D(p_i)\chi_q(p_i) = \chi_m(p_i) = \chi_m^\circ$ --- главный характер


(т.\,е. $\delta_{\chi}=1$) тогда и только тогда, когда


$\chi_D(p_i)=\ov \chi_q(p_i)$ и $\chi_q(p_i)=\chi_d(p_i)$ для всех


$(p_i, m)=1$. Поэтому, применяя лемму 2.2 при $Q'=D$, $Q''=q$, $Q=m$, из


(5.5) получим


$$


A_i=\frac1{\varphi(q)}


\sum_{\chi_D, \chi_q}\sum_{j=1}^q {\ov \chi_{D}}(l_i){\ov \chi_q}(j)


e\bigg(\frac{a j}{q}\bigg)=


%


\frac1{\varphi(q)}


\sum_{\chi_d}\sum_{j=1}^q \chi_d(l_i){\ov \chi_d}(j)


e\bigg(\frac{a j}{q}\bigg)=\frac{\varphi(d)}{\varphi(q)}


\sideset{}{'}\sum\begin{Sb}j=1\\ j\equiv l_i(\Mod d)\end{Sb}^q


e\bigg(\frac{a j}{q}\bigg).


$$


Отсюда, используя лемму 2.3, находим


\begin{equation}


A_i=R(D)\frac{\mu(qd^{-1})}{\varphi(qd^{-1})}


e\bigg(\frac{a N_1l_i}{q}\bigg).


\end{equation}


Из (3.3), (5.9) и (5.10) следует утверждение леммы 5.4


\end{pf}





\begin{lems}


Справедлива оценка


$$


\sum_{q\le P}\sideset{}{'}\sum_{a=1}^{q}\int_{M(q,a)}


|S(N,\alpha) - V(N,\alpha,q,a)|^2d\alpha\ll N^3P^3D^2\exp(-2c_6\sqrt{\ln N}).


$$


\end{lems}





{\bf Доказательство.} Утверждение леммы следует из леммы 5.4, если учесть


в силу (3.3) и (3.10)


$$


V_i(N,\alpha,q,a)\ll \varphi^{-1}(qd^{-1})|\mathrm g_1^{(i)}(N,\eta)|\ll


N\varphi^{-1}(qd^{-1}).


$$


Действительно, из леммы 5.4 следует


\begin{multline*}


S(N,\alpha) - V(N,\alpha,q,a)\ll(|V_1(N,\alpha,q,a)|+


|V_2(N,\alpha,q,a)|)NPq^{-1}\exp(-c_6\sqrt{\ln N})+\\


%


+(NPq^{-1}\exp(-c_6\sqrt{\ln N}))^2\ll N^2P\varphi^{-1}


(qd^{-1})


q^{-1}\exp(-c_6\sqrt{\ln N}) +N^2P^2q^{-1}\exp(-2c_6\sqrt{\ln N}).


\end{multline*} 


Отсюда


\begin{multline*}


\sum_{q\le P}\,\sideset{}{'}\sum_{a=1}^{q}\int_{M(q,a)}


|S(N,\alpha) - V(N,\alpha,q,a)|^2d\alpha\ll \\


%


\ll \frac{N^4P^2}{\exp(2c_6\sqrt{\ln N})}


\sum_{q\le P}\frac{\varphi(q)}{q^3Q\varphi^2(qd^{-1})}+


\frac{N^4P^4}{\exp(4c_6\sqrt{\ln N})}\sum_{q\le P}\varphi(q)q^{-5}Q^{-1}\ll\\


%


\ll N^3P^3D^2\exp (-2c_6\sqrt{\ln N}).\quad\square


\end{multline*}





\begin{lems}


Справедливо соотношение


$$


\int_{Q^{-1}}^{1+Q^{-1}}


\bigg|S(N,\alpha) - \sum_{q=1}^{\infty}


\sideset{}{'}\sum_{a=1}^q V(N,\alpha, q, a)\bigg|^2 d\alpha \ll


N^3D^2P^{-1}\ln^{12} N.


$$


\end{lems}





\begin{pf}


В силу леммы 5.3 оцениваемый интеграл можно представить в виде


\begin{equation}


\int_{Q^{-1}}^{1+Q^{-1}}


\bigg|S(N,\alpha) - \sum_{q\le P}


\sideset{}{'}\sum_{a=1}^q V(N,\alpha, q, a)\bigg|^2 d\alpha +


O\bigg(\frac{N^3D^4\ln^4 N}{P^2}\bigg).


\end{equation}





Пусть $M=[Q^{-1},1+Q^{-1}]\setminus T$. Тогда интеграл по $[Q^{-1},1+Q^{-1}]$


в (5.11) можно записать как сумму интегралов по $M$ и по $T$, которые обозначим


через $K_1$ и $K_2$ соответственно.





$K_2$ оценим при помощи леммы 4.1 


\begin{equation} 


K_2\ll\int_{T}|S(N,\alpha)|^2d\alpha +\int_{T}\bigg|


\sum_{k\le P}\sideset{}{'}\sum_{h=1}^kV(N,\alpha,k,h)\bigg|^2 d\alpha\ll


N^3D^2P^{-1}\ln^{12} N.


\end{equation}





Для оценки интеграла $K_1$ используем леммы 5.1 и 5.5. Это дает


\begin{multline} 


K_1\ll\int_{M}\bigg|S(N,\alpha)-


\underset{\alpha\in M(k,h)}{\sum_{k\le P}\sideset{}{'}\sum_{h=1}^k}{}


V(N,\alpha, k, h)\bigg|^2 d\alpha +\\


+\int_M


\bigg|\underset{\alpha\notin M(k,h)}


{\sum_{k\le P}\sideset{}{'}\sum_{h=1}^k}{}


V(N,\alpha, k, h)\bigg|^2 d\alpha\ll N^3P^3D^2\exp(-2c_6\sqrt{\ln N}).


\end{multline}


Теперь из (5.11)--(5.13) следует утверждение леммы 5.6.


\end{pf}





Б. В случае $P=P_1$ получится $E=1$ и существует исключительный 


вещественный (нуль $\beta$) характер $\wt \chi_r$, где


$r\le P_1^{1/4}=P_2$.


Пусть $\tau(\chi_k)$ --- сумма Гаусса, т.\,е.


$$


\tau(\chi_k)=\sum_{n=1}^k \chi_k(n)e(n/k).


$$


Из теоремы Пейджа и Зигеля о нулях $L$-функции Дирихле [6] следует, что


$1-c_7\ln^{-1/2}N<\beta<1-c(\varepsilon)r^{-\varepsilon}$, т.\,е.


$r\gg (\ln N)^{1/2\varepsilon}$.


Положим $\varepsilon= \varepsilon_1=(2A+6)^{-1}$, тогда


$r>(\ln N)^{A+2}$.





Известно, что ведущие модули вещественного примитивного характера могут


равняться только 4, 8 и $p>3$ --- простое число. Так как


$m=qDd^{-1}$, $d=(q,D)$, $D=p^\nu$, $D\ll \ln^a X$,


$(qd^{-1}, D)=1$, $qd^{-1}$


бесквадратный, то отсюда следует, что $r\setminus m$ равносильно тому,


что $r\setminus qd^{-1}$, т.\,е. $r\setminus q$. 


\medskip





{\bf Лемма 5.$\wt{\bold 1}$.} $\mathrm a)$


{\it если $r\not\setminus q$, то}


$$


\sum_{q\le P}\sideset{}{'}\sum_{a=1}^q \int_{M(q,a)}


|S(N,\alpha) - V(N,\alpha, q, a)|^2 d\alpha\ll 


\frac{N^3P^3D^2}{\exp(2c_6\sqrt{\ln N})};


$$





$\mathrm b)$ {\it если $r\setminus q$, то


$$


\sum_{q\le P}\sideset{}{'}\sum_{a=1}^q \int_{M(q,a)}


\bigg|S(N,\alpha) - \prod_{i=1}^2(V_i-\wt V_i)\bigg|^2 d\alpha\ll 


\frac{N^3P^3\ln P}{\exp(2c_6\sqrt{\ln N})},


$$


где}


$$


\wt V_i =A_i(\chi_q, \chi_{D}, q, a)\mathrm g_\beta^{(i)}(N,\eta).


$$





Доказывается так же, как и лемма 5.5. При этом в доказательстве


утверждения a) используется лемма 2.1 a), а в доказательстве утверждения


b) --- лемма 2.1 b).


\medskip





{\bf Лемма 5.$\wt{\bold 2}$} ([3], лемма 5.2).


{\it Если $\chi_{k}$ --- 


характер по $\Mod k$, индуцированный примитивным характером 


$\chi_{k}^*$, то


$r\setminus k$ и


$\tau (\chi_{k})=\mu (k/r)\chi_{r}^*(k/r)\tau(\chi_r^*)$,


$|\tau (\chi_r^*)|^2=r$.}


\medskip





Из равенства (5.5) следует





{\bf Лемма 5.$\wt{\bold 3}$.} 


{\it Если существует исключительный характер по модулю $m=qDd^{-1}$, то}


$$


A_i(\chi_{D}, \chi_q, q, a)=\frac1{\varphi(q)}\wt \chi_m


(al_i)\sideset{}{'}\sum_{h=1}^q\wt \chi_m(h)e(h/q).


$$





Далее для $a\le q\le P$, $(a,q)=1$, $r\setminus m$,


$\alpha\in M(q,a)$ положим 


$W(N,\alpha)=\wt V_1\wt V_2 -V_1\wt V_2 +\wt V_1V_2$, в других


случаях $W(N,\alpha)=0$.


\medskip





{\bf Лемма 5.$\wt{\bold 4}$.} {\it Справедлива оценка}


$$


\int_{Q^{-1}}^{1+Q^{-1}}


\bigg|S - \sum_{q=1}^{\infty}\sideset{}{'}\sum_{a=1}^q 


V(N,\alpha, q, a)-W(N,\alpha)\bigg|^2 d\alpha \ll


N^3D^2P^{-1}\ln^{12} N.


$$





\begin{pf}


В силу леммы 5.3 оцениваемый интеграл представим в виде


\begin{equation}


\int_{Q^{-1}}^{1+Q^{-1}}


\bigg|S - \sum_{q\le P}\sideset{}{'}\sum_{a=1}^q 


V(N,\alpha, q, a)-W(N,\alpha)\bigg|^2 d\alpha +


O(N^3P^{-2}(D\ln N)^4).


\end{equation}





Интеграл по $[Q^{-1}, 1+Q^{-1}]$ в (5.14) представим суммой интегралов


по $M$ и $T$. Тогда учитывая, что $W(N,\alpha)=0$ для $\alpha\in T$, и


используя леммы 4.1, 5.1 и $5.\wt 1$, получим утверждение леммы


(подробности --- в [2]). 


\end{pf}





\section{Исследование функции $W(N,\alpha)$} 





Положим


\begin{equation}


D(N, h)=\int_{Q^{-1}}^{1+Q^{-1}}W(N,\alpha)e(-h\alpha)d\alpha.


\end{equation} 





\begin{lems}


Имеет место оценка


$$


\sum_{2<n\le 2N}|\varphi^2(D)R(N,n)-J(N,n)\sigma(n)-D(N,n)|^2


\ll N^3D^3P^{-1}\ln^{12} N,


$$


где $R(N,n)$, $J(N,n)$ и $\sigma(n)=\sigma(0,n)$ определены в виде сумм 


$(3.6)$, $(3.8)$ и $(5.2)$.


\end{lems}





\begin{pf}


Пусть


$$


F(N,h)=


\begin{cases}


\varphi^2(D)R(N,h)-J(N,h)-D(N,h),\ \; 0<h\le 2N;\\


D(N,h)\; \text{ в других случаях}.


\end{cases}


$$


Тогда в силу (6.1), (5.1) и (3.5) функция $F(N,h)$


является коэффициентом Фурье функции


$$


S - \sum_{q=1}^{\infty}\sideset{}{'}\sum_{a=1}^q 


V(N,\alpha, q, a)-W(N,\alpha).


$$


Применяя неравенство Бесселя и лемму $5.\wt 4$, получим утверждение


леммы.


\end{pf}





Пусть $r\setminus m$ и


\begin{gather}


I_1^{(i)}=I_1^{(i)}(n,q,D,a)=\int_{M(q,a)}V_i\wt V_i


e(-n\alpha)d\alpha,\quad i=1,2;\\


%


I_2^{(i)}=I_2^{(i)}(n,q,D,a)=\int_{M(q,a)}\wt V


e(-n\alpha)d\alpha,\quad \wt V=\wt V_1\wt V_2; \\


%


B_1^{(i)}=B_1^{(i)}(n,q,D,a)=\mu(qd^{-1})\varphi^{-1}(qd^{-1})


e(aq^{-1}(N_1l_i-n))A_i,\quad i=1,2;\\


%


B_2=B_2(n,q,D,a)=A_1\cdot A_2e(-aq^{-1}n),


\end{gather}


тогда


\begin{equation}


I_1^{(i)}=B_1^{(i)}\int_{-1/qQ}^{1/qQ}\mathrm g_1^{(i)}(N,\eta)


\mathrm g_{\beta}^{(i)}


(N,\eta) e(-n\eta)d\eta,\quad i=1,2,


\end{equation}


и


\begin{equation}


I_2=B_2\int_{-1/qQ}^{1/qQ}\mathrm g_{\beta}^{(1)}(N,\eta)


\mathrm g_{\beta}^{(2)}


(N,\eta)e(-n\eta)d\eta.


\end{equation}





Теперь, если $q\le P$, $r\setminus qd^{-1}$ и $k=qr^{-1}$, то из (6.4)


и леммы $5.\wt 3$ следует оценка


$$


\sideset{}{'}\sum_{a=1}^qB_1^{(i)}(n,q,D,a)\ll \frac{r^{1/2}}{\varphi^2(q)}


\varphi(d)\varphi(r)\wt \chi_r(k)^2|C_k(N_1l_i-n)|\mu^2(k).


$$





Аналогично из (6.5) и леммы $5.\wt 3$ получим


\begin{equation}


\sideset{}{'}\sum_{a=1}^qB_2(n,q,D,a)\ll \frac{1}{\varphi^2(q)}


\wt \chi_r(l_1l_2)\mu^2(k)\wt \chi_r^2(k)\tau(\wt \chi_r)^2


C_q(-n),


\end{equation}


где 


\begin{equation}


C_q(m)=\sideset{}{'}\sum_{a=1}^q e\bigg(\frac{am}q\bigg)=


\frac{\mu(q_1)}{\varphi(q_1)}


\varphi(q),\quad q_1=\frac{q}{(q,|m|)}.


\end{equation}





\begin{lems}[{\series{m}\shape{n}\selectfont [2], лемма 9.4}]


Пусть $L_i=|N_1l_i-n|$, тогда


$$


\sum_{k\le Pr^{-1}}\mu^2(k)\varphi^{-2}|C_k(N_1l_i-n)|\ll


L_i\varphi^{-1}(L_i).


$$


\end{lems}





\begin{lems} Если $r\setminus m$, то


$$


\sum_{q\le P}\sideset{}{'}\sum_{a=1}^q I_1^{(i)}(n,q,D,a)\ll


Nr^{1/2}\varphi(d) L_i(\varphi(r)\varphi(L_i))^{-1}.


$$


\end{lems}





\begin{pf}


Используя неравенство Шварца и (3.2) в (6.6), а затем


применяя лемму 6.2, получим утверждение леммы.


\end{pf}





Пусть


\begin{equation}


J_\beta(N,n)=\int_{-1/2}^{1/2}\mathrm g_{\beta}^{(1)}\mathrm g_{\beta}^{(2)}


e(-n\eta)d\eta=


\sum\begin{Sb}n=n_1+n_2\\ 2\le n_1,n_2\le N\end{Sb}(n_1n_2)^{\beta-1}.


\end{equation}


Тогда


\begin{equation}


\int_{-1/qQ}^{1/qQ}\mathrm g_{\beta}^{(1)}\mathrm g_{\beta}^{(2)}


e(-n\eta)d\eta=J_\beta(N,n)+O\bigg(\frac{qN}{P}\bigg).


\end{equation}


Далее положим


\begin{equation}


G(N,n)=\sum\begin{Sb}q\le P\\ r\setminus qd^{-1}\end{Sb}


\sideset{}{'}\sum_{a=1}^q B_2(n,q,D,a).


\end{equation}





\begin{lems}


Если $n\le N$, то


$$


\sum\begin{Sb}q\le P\\ r\setminus q\end{Sb}


\sideset{}{'}\sum_{a=1}^q I_2(n,q,D,a)=J_\beta(N,n) G(N,n)


+O\bigg(\frac{qN}{P}\tau(n)(\ln\ln N)^4\ln^{1/2} N\bigg).


$$


\end{lems}





\begin{pf}


Согласно (6.7) и (6.11) из (6.8), (6.9) и (6.12)


получим


\begin{equation}


\sum\begin{Sb}q\le P\\ r\setminus qd^{-1}\end{Sb}


\sideset{}{'}\sum_{a=1}^q I_2(n,q,D,a)-J_\beta(N,n) G(N,n)\ll


\frac{Nr^2}{P\varphi(r)}\sum\nolimits_1,


\end{equation}


где 


\begin{equation}


\sum\nolimits_1 =\sum\begin{Sb}k\le Pr^{-1}\\ (k,r)=1\end{Sb}


\mu^2(k)\mu^2\bigg(\frac{k}{(k,n)}\bigg)


\frac{k}{\varphi(k)}\varphi^{-1}\bigg(\frac{k}{(k,n)}\bigg)\ll


(\ln\ln N)^3\tau(n)\ln^{1/2}N.


\end{equation}


Так как $r\varphi^{-1}(r)\ll\ln\ln N$, то из (6.13) и (6.14)


следует утверждение леммы.


\end{pf}





\begin{lems}


Если $n\le 2N$, то


$D(N,n)-J_\beta(N,n) G(N,n)\ll


NP^{-1}r\tau(n)(\ln\ln N)^4\ln^{1/2}N+ Nr^{-1/2}\varphi(d)(\ln\ln N)^2$.


\end{lems}





{\bf Доказательство.} В силу (6.1)--(6.3) имеем


$$


D(N,n)=\sum_{q\le P}\sideset{}{'}\sum_{a=1}^q 


(I_2-I_1^{(1)}-I_1^{(2)}).


$$


Далее, применяя леммы 6.3 и 6.4, будем иметь


$$


D(N,n)-J_\beta(N,n) G(N,n)\ll


NP^{-1}\tau(n)r(\ln\ln N)^4\ln^{1/2}N+


Nr^{1/2}\varphi(d)\varphi^{-1}(r)\ln\ln N.


$$


Отсюда следует утверждение леммы 6.5, если учесть, что 


$$


\varphi(r)\gg r(\ln\ln r)^{-1},\quad


\text{где}


\quad P\ge r>\ln^{A+2}N.\quad\square


$$





\section{Доказательство теоремы 1}





A. Случай $P=P_2$.


Из (3.5) и (5.1) при $y=0$, используя тождество Парсеваля и лемму 5.6,


находим


\begin{equation}


\sum_{n\le 2N}(\varphi^2(D)R(N,n)-J(N,n)\sigma(n))^2 \ll\frac{(ND)^2}{P}


\varphi(D)\ln^{12}N.


\end{equation}


Здесь для четного $n$


\begin{multline*}


\sigma(n)=\sigma(0,n)=D\sum_{k=1}^{\infty}\mu(k)\varphi^{-2}(k)


\sum_{t\setminus n,\; (t,D)=1}\mu^2(t)\varphi^{-1}(t)= \\


%


=\lambda D\prod\limits_{p>2}\frac{p(p-2)}{(p-1)^2}


\prod\begin{Sb}p\setminus D\\ p>2\end{Sb}\frac{(p-1)^2}{p(p-2)}


\prod\begin{Sb}p\setminus n, p\not\setminus D\\ p>2\end{Sb}\frac{p-1}{p-2},


\end{multline*}


где $\lambda = 1$, если $D$ четное; при $D$ нечетном $\lambda = 2$


(более подробно относительно $\sigma(n)$ см. \S\,7 из [1]).


Отсюда


\begin{equation}


\sigma(n)\gg D.


\end{equation}


Из (7.1) следует, что для всех $n\le 2N$, исключая не более чем


\begin{equation}


c_8 NP^{-1/3}\varphi(D)\ln^4 N


\end{equation} 


значений из них, справедлива формула


\begin{equation}


R(N,n)=\frac{J(N,n)\sigma(n)}{\varphi^2(D)}+


O\bigg(\frac{N(\ln N)^4\ln\ln D}{\varphi(D)\exp(c_6\sqrt{\ln N}/600)}\bigg).


\end{equation}


Согласно (7.2) и (3.9) при $N/2<n\le N$ имеем


$$


J(N,n)\sigma(n)\gg DN.


$$


Тогда из (7.4) находим


\begin{equation}


R(N,n)\gg DN\varphi^{-2}(D)\bigg(1-P^{-1/3}\ln^4N\bigg)\gg 


N\varphi^{-1}(D).


\end{equation}


Из (7.3) и (7.5) сразу следует утверждение теоремы 1 с $X=N$.


Из (7.5) следует, что


\begin{equation}


R(n)\gg n(\varphi(D)\ln^2n)^{-1}\big(1-(\ln n)^2\exp(-c_9\sqrt{\ln n})\big).


\end{equation}


\medskip





Б. Случай $P=P_1$. В силу (6.13) и (6.8) имеем


\begin{equation}


|G(N,n)|\le\wt\sigma(n)+O(P^{-1} r d\tau(n)(\ln\ln N)^4),


\end{equation}


где


$$


\wt\sigma(n)=\prod\begin{Sb}p\setminus rd\\


p\not\setminus n\end{Sb}\frac1{p-1}


\prod_{p\not\setminus nrd}\frac{p(p-2)}{(p-1)^2}


\prod_{p\setminus nrd}\frac{p}{p-1}\le \sigma(n)


$$


(подробные выкладки см. в [1] и [2]).





Учитывая $J_\beta(N,n)\ll N^\beta$ (см. (6.10)), из леммы 6.5 и неравенства


(7.7) получим для четного $n\le 2N$


\begin{equation}


|D(N,n)|\le J_\beta(N,n)\sigma(n)+


O\big(P^{-1} d r\tau(n)(\ln\ln N)^4\ln^{1/2}N\big)+


O\big(Nr^{-1/2}\varphi(D)(\ln\ln N)^2\big).


\end{equation}





\begin{lems}


Существует такая положительная постоянная $c_{10}$, что


$$


J(N,n)-J_\beta(N,n) >c_{10}r^{-\varepsilon_1}J(N,n),\quad 


\varepsilon_1= (2A+6)^{-1}.


$$


\end{lems}





\begin{pf}


В силу теоремы Зигеля $1/2\le\beta<1-c(\varepsilon_1)r^{-\varepsilon_1}$.


Поэтому при $n_1n_2>1$ будем иметь


$$


\big(1-(n_1n_2)^{\beta-1}\big)>1-


\exp\big(-c(\varepsilon_1)r^{-\varepsilon_1}\ln(n_1n_2)\big)


>c_{10}r^{-\varepsilon_1}.


$$


Теперь, используя последнюю оценку, из (6.10) и (3.8) получим утверждение


леммы.


\end{pf}





Из (7.8) и леммы 7.1 следует, что если $n\le 2N$ и $n\equiv 0(\Mod 2)$, то


\begin{multline}


|J(N,n)\sigma(n)+D(N,n)|>c_{10}r^{-\varepsilon_1}J(N,n)\sigma(n)+ \\


%


+ O\bigg(\frac{N}{P}r d \tau(n)(\ln\ln N)^4\ln^{1/2}N\bigg)+


O\bigg(\frac{N}{r^{1/2}}\varphi(D)(\ln\ln N)^2\bigg).


\end{multline}


Согласно лемме 6.1


\begin{equation}


\varphi^2(D)R(N,n)- J(N,n)\sigma(n)-D(N,n)\ll NP^{-1/3}D\varphi(D)


\end{equation}


для всех $n\le 2N$, за исключением не более чем


\begin{equation}


c_{11} \varphi^{-1}(D)NP^{-1/3}\ln^{12}N


\end{equation}


значений из них.





\begin{lems}


Для всех четных чисел $n$, $N/2<n\le N$, 


за исключением не более чем


$c_{12}N \varphi^{-1}(D)P^{-1/3}\ln N$ значений из них, справедлива


оценка


$$


|J(N,n)\sigma(n)+D(N,n)|\gg NP^{-\varepsilon_1/4}D.


$$  


\end{lems}





\begin{pf}


Известно, что $\sum\limits_{n\le N}\tau(n)\ll N\ln N$.


Отсюда $\tau(n)\ll P^{1/3}\varphi(D)$ для всех $n$, $n\le N$,


за исключением не более чем


$c_{12}N P^{-1/3}\varphi^{-1}(D)\ln N$ значений $n$.





Следовательно, согласно (7.9)


\begin{multline*}


|J(N,n)\sigma(n)+D(N,n)|>c_{10}J(N,n)\sigma(n)+


O(Nr^{-1}P^{-5/12}\varphi(D) d


(\ln\ln N)^4\ln^{1/2}N)+ \\


%


+O(Nr^{-1/2}(\ln\ln N)^2\varphi(d))\gg r^{-\varepsilon_1}ND


(1-c_{13}(\ln\ln N)^2(\ln N)^{-(A+2)(\frac12-\varepsilon_1)})


\gg r^{-\varepsilon_1}ND.


\end{multline*}


Так как $r\le P$, то отсюда следует утверждение леммы.


\end{pf}





Поскольку $R(N,n)$ не отрицательное, то


$$


R(N,n)\ge\varphi^{-2}(D)\{|J(N,n)\sigma(n)+D(N,n)|-


|\varphi^2(D)R(N,n)-J(N,n)\sigma(n)-D(N,n)|\}.


$$         





Следовательно, учитывая $P=P_1$, из (7.10) и леммы 7.2 получим


$$


R(N,n)>\frac{ND}{P^{\varepsilon_1/4}}\big(c_{14}-c_{15}\varphi(D)


P^{-\frac13+\frac{\varepsilon_1}{4}}\big)> NDP^{-\frac{1}{20}}.


$$


Из последнего соотношения следует


\begin{equation}


R(n)>\frac{n\exp(-c_{14}\sqrt{\ln n})}{\varphi(D)\ln^2n}\bigg(


1-\frac{\ln^A n}{\exp(3c_{14}\sqrt{\ln n})}\bigg).


\end{equation}


Из (7.6) и (7.12), учитывая (7.3) и (7.11), получим


утверждение теоремы 1.$\quad\square$





В заключение автор выражает благодарность за поддержку и полезные советы


профессорам А.Ф.\,Лаврику и М.И.\,Исраилову.





\begin{thebibliography}{99}





\bibitem{1}


Лаврик А.Ф. {\em К бинарным проблемам аддитивной теории простых


чисел в связи с методом тригонометрических сумм И.М.\,Виноградова} //


Вестн. Ленингр. ун-та. Сер. матем. -- 1961. -- Вып.\,13. -- С.\,11--27.


\bibitem{2}


Vaughan R.C. {\em On Goldbach's problem} // Acta arithm. -- 1972. --


V.\,21. -- \No\,1. -- P.\,21--48.


\bibitem{3}


Montgomery H.L., Vaughan R.C. {\em The exceptional set in


Goldbach's problem} // Acta arithm. -- 1975. -- V.\,27. -- P.\,353--370.


\bibitem{4}


Аллаков И.А. {\em Некоторые оценки снизу для числа представлений


гольдбаховых чисел} //


Вопр. вычисл. и прикл. матем. -- Ташкент. -- 1985. -- \No\,77.


-- C.\,37--41.


\bibitem{5}


Хамзаев Э. {\em О представлении натуральных чисел в виде суммы


простого и квадрата целого из арифметической прогрессии} //


Узб. матем. журн. -- 1991. -- \No 3. -- C.\,64--75.


\bibitem{6}


Дэвенпорт Г. {\em Мультипликативная теория чисел}. -- М.: Наука,


1971. -- 199~с.


\bibitem{7}


Rademacher H. {\em \"Uber eine Erweiterung des


Goldbachen Problems} // Math. Zeit. -- 1925. -- Bd.\,25. -- S.\,627--657.


\bibitem{8}


Аллаков И.А. {\em Об исключительном множестве в бинарной проблеме


Гольдбаха}. -- Ред. журн. ``Изв. АН УзССР. Сер. физ.-матем. наук''. --


Ташкент, 1981. -- 76 с. -- Деп. в ВИНИТИ 11.11.81, \No\,5087--82.


\bibitem{9}


Прахар К. {\em Распределение простых чисел}. -- М.: Мир, 1967. -- 511 с.


\bibitem{10}


Исраилов М.И., Аллаков И.А. {\em Об оценке тригонометрических сумм по 


простым числам арифметических прогрессий} //


ДАН УзССР. -- 1982. -- \No\,4. -- C.\,5--6.





\end{thebibliography}





\vskip 2cm





\postup{\it Термезский государственный университет}{\it Поступили}


\postup{}{\it первый вариант $17.02.1997$}


\postup{}{\it окончательный вариант $29.02.2000$}





\label{end}


\end{document}


